\chapter{Dying at Home}

\section*{Definition and Overview}
Dying at home means spending the final days or weeks of life in the familiar surroundings of one's own home, supported by family and professional carers, rather than in hospital or a hospice. Many people express a wish to die at home, and with good planning and support this can be achieved for a significant proportion of people with life-limiting illnesses. This chapter explains what is involved in caring for a dying person at home, the services that support home death, and the practical considerations families should plan for in advance.

\section*{Epidemiology}
A large proportion of Australians say they would prefer to die at home, but only a minority of deaths currently occur there, with most happening in hospitals. The likelihood of dying at home depends on the illness, the availability of family carers, access to palliative care services, and the wishes and preparedness of the family. With multidisciplinary support, home death is feasible for many people, particularly those with cancer, and rates have increased with the growth of community palliative care.

\section*{Aetiology and Pathophysiology}
The physical process of dying at home is the same as dying anywhere; what differs is the environment and the support available. As death approaches, the body progressively shuts down, with reduced energy, appetite, circulation, and consciousness. Understanding these changes helps families recognise what is normal and respond with calm and confidence.

\section*{Clinical Features}
Families caring for a dying person at home should be familiar with the changes that occur in the final days.
\begin{itemize}
    \item Increasing drowsiness and unresponsiveness, with hearing often remaining intact
    \item Reduced interest in food and drink
    \item Irregular breathing and pauses between breaths
    \item A rattling sound from the throat due to secretions
    \item Cool, mottled hands and feet
    \item Confusion, restlessness, or agitation at times
    \item Reduced urine output
    \item Changing pain levels, requiring adjustment of medicines
\end{itemize}

\section*{Differential Diagnosis}
Not every change is part of the natural process of dying, and some are treatable.
\begin{itemize}
    \item Reversible infections, such as urinary or chest infections, which may cause confusion or fever
    \item Medication side effects, including drowsiness and nausea
    \item Uncontrolled pain or breathlessness requiring changes to medication
    \item Urinary retention or constipation causing distress and agitation
    \item Anxiety and distress, which respond to reassurance and appropriate medication
\end{itemize}

\section*{When to Seek Medical Care}
Families should have clear instructions on who to call and when. The community palliative care team, GP, or after-hours services should be contacted for uncontrolled symptoms, concerns about medication, or any change that is distressing.

\textbf{Call 000 (or local emergency services) for:}
\begin{itemize}
    \item Sudden severe breathlessness, chest pain, or collapse
    \item Sudden severe bleeding
    \item Seizures or fits
    \item The person cannot be roused and this is unexpected
    \item Any emergency situation where immediate assistance is needed
\end{itemize}

\section*{Diagnosis and Investigations}
Assessment at home is focused on comfort rather than investigation.
\begin{itemize}
    \item \textbf{Regular nursing review:} nurses assess pain, breathing, comfort, and family confidence
    \item \textbf{Symptom assessment:} a structured approach to pain, nausea, agitation, and breathing difficulties
    \item \textbf{Medication planning:} ensuring appropriate medicines are available at home, including those for breakthrough symptoms
    \item \textbf{Advance care planning:} confirming the person's wishes about resuscitation, transfer to hospital, and treatment
\end{itemize}

\section*{Management}
Successful home care depends on planning, equipment, and the involvement of a professional team.
\begin{itemize}
    \item \textbf{Palliative care team:} community nurses and doctors provide regular visits, symptom management, and 24-hour advice
    \item \textbf{Equipment:} hospital bed, commode, pressure-relieving mattress, and other aids arranged through palliative care services
    \item \textbf{Medicines:} a "just in case" supply of medicines for pain, agitation, nausea, and secretions, with clear instructions
    \item \textbf{Family education:} training in medication administration, positioning, mouth care, and what to do at the time of death
    \item \textbf{Practical support:} help from volunteers, respite care, and allied health services
    \item \textbf{After-death care:} clear guidance on who to call and what happens after the person dies at home
\end{itemize}

\section*{Complications}
\begin{itemize}
    \item Carer exhaustion and burnout, which can jeopardise the ability to continue home care
    \item Uncontrolled symptoms if services are not well coordinated
    \item Family disagreement about care decisions
    \item A transfer to hospital at the very end if plans are not clear
    \item The emotional burden of caring, which requires its own support
\end{itemize}

\section*{Prognosis and Outcomes}
With good planning and support, dying at home can be a peaceful and meaningful experience for both the person and their family. Community palliative care achieves good symptom control in the majority of cases, and families who have cared for a loved one at home often report that, despite the difficulty, they value the time spent together. Clear advance care planning greatly increases the chance that death occurs at home as wished.

\section*{Prevention and Self-Care}
\begin{itemize}
    \item Discuss the wish to die at home early and document it in an advance care directive
    \item Engage palliative care services early, rather than waiting until the final days
    \item Organise equipment and medicines before they are needed
    \item Plan a roster of family and friends so no single carer becomes exhausted
    \item Ask the team what will happen at the moment of death, so there are no surprises
    \item Accept offers of help with meals, cleaning, and errands
    \item Look after your own health, rest, and emotional support
\end{itemize}

\section*{Sources and Further Reading}
The following reputable sources provide further information for healthcare students and patients seeking reliable, current advice about dying at home.
\begin{itemize}
    \item Palliative Care Australia. \textit{Dying at home: planning and support.} https://palliativecare.org.au/
    \item NHS. \textit{Caring for someone at the end of life at home.} https://www.nhs.uk/conditions/end-of-life-care/
    \item CareSearch. \textit{Dying at home: information for families.} https://www.caresearch.com.au/
    \item Advance Care Planning Australia. \textit{Planning for end-of-life care.} https://www.advancecareplanning.org.au/
    \item Mayo Clinic. \textit{End-of-life care at home.} https://www.mayoclinic.org/patient-visitor-guide/
    \item MedlinePlus. \textit{End of life care: providing comfort at home.} https://medlineplus.gov/endoflifeissues.html
\end{itemize}
